\documentclass[12pt]{article}
\usepackage[utf8]{inputenc}
\usepackage[T1]{fontenc}
\usepackage{hyperref}
\title{Tirahi Language Resources for K-12 Curriculum Development}
\author{Compiled by Claude Sonnet 5 (claude-sonnet-5) for Andr\'as Kornai}
\date{}
\begin{document}
\maketitle
\section{Introduction}
Tirahi (ISO 639-3 tra) is a nearly extinct Indo-Aryan (Dardic) language once spoken in a few villages southeast of Jalalabad, Nangarhar Province, Afghanistan. Sources disagree on the number of speakers. One lists about 100, while another lists zero and calls the language dormant. The remaining speakers are older adults who also speak Pashto. The language is unwritten. The search found no school materials.
\section{Primary Educational Resources}
No primers, textbooks or school programmes were found in the search results.
\section{Linguistic and Reference Materials}
The search results show only encyclopedic overviews. Morgenstierne's 1934 ``Notes on Tirahi'' is mentioned there as the key early description, but no link to the work itself was found.
\begin{itemize}
\item Wikipedia, Tirahi language: \url{https://en.wikipedia.org/wiki/Tirahi_language}
\item Endangered Languages Project entry: \url{https://www.endangeredlanguages.com/lang/3509?hl=zh-cn}
\item Verbix Endangered Languages entry: \url{https://docs.verbix.com/EndangeredLanguages/Tirahi}
\end{itemize}
\section{Religious and Community Resources}
Joshua Project has a people-group profile, which says that most younger Tirahi have shifted to Pashto. \url{https://joshuaproject.net/people_groups/15500/AF}
\section{Digital and Audio Resources}
No audio, video or corpus resources were found.
\section{Curriculum Development Context}
Very little exists for Tirahi. Any heritage or revitalization curriculum would first need documentation work with the remaining elderly speakers, if there are any. Existing sources do not support a standard curriculum.
\end{document}